% GENERATED FILE. Edit canonical lessons, then run npm run generate:books.
% canonical-source-hash: fnv1a64:d969a8d70e270eb2
% canonical-lessons: ZH-C31-ta-she, ZH-C31-ta-he, ZH-W31-ta, ZH-C31-tamen, ZH-W31-men, ZH-C31-women, ZH-C31-nimen

\chapter{He, She and We}
\label{ch:zh-he-she-and-we}
\hlchaptermodality{31}

\begin{chapteropening}
\textbf{By the end of this chapter:} \emph{I can say he, she, we, you all and they.}

Complete the last lesson of chapter 31: I can say he, she, we, you all and they.
\end{chapteropening}

\section[tā]{tā — she; her}

\label{lesson:ZH-C31-ta-she}



\begin{quote}
Before the new one: say the Mandarin for tired, then the Mandarin for not bad.
\end{quote}

\subsection*{You'll want to know: tā}
\emph{tā} — \textquotedblleft{}she; her\textquotedblright{}.

\begin{quote}
\emph{tā shì wǒ de tóngxué} — She's a classmate of mine.
\end{quote}

Its characters come later; for now it is a word for the ear.

\begin{cousinweb}
\emph{tā} sounds exactly like the word for he. No dictionary we use records an older form, so its history is not settled.
\end{cousinweb}

\begin{grammarlens}[title={What you've built}]
One more word about people.
\end{grammarlens}

\subsection*{Your turn}
\begin{itemize}
\raggedright
  \item \emph{Say it:} \emph{tā}
  \item \emph{Say it:} \emph{tā}, once more
  \item \emph{Say it:} \emph{tā shì wǒ de tóngxué} — She's a classmate of mine.
  \item \emph{Recall:} say the Mandarin for tired, then the Mandarin for not bad, then say \emph{tā} again
\end{itemize}

\begin{tcolorbox}[breakable,colback=teal!4,colframe=teal!35!black,title={Before you move on}]
What does tā mean? (\textquotedblleft{}She; her\textquotedblright{}.) Say it once more.
\end{tcolorbox}
\section[\texorpdfstring{\zh{他} (tā)}{他 (tā)}]{\zh{他} (tā) — he; him}

\label{lesson:ZH-C31-ta-he}



\begin{quote}
Before the new one: say the Mandarin for not bad, then the Mandarin for she.
\end{quote}

\subsection*{You'll want to know: \zh{他}}
\textbf{\zh{他}} — \emph{tā} — \textquotedblleft{}he; him\textquotedblright{}.

\begin{quote}
\textbf{\zh{他是我的儿子}} \emph{tā shì wǒ de érzi} — He's my son.
\end{quote}

The new character, in pieces: \zh{他} — the person piece from \zh{你} on the left, and on the right a three-stroke shape read yě. The next lesson puts it in your hand.

\begin{cousinweb}
\textbf{\zh{他}} (\emph{tā}) is inherited from Middle Chinese \emph{tha} (Guangyun, in Baxter's transcription). The dictionary glosses it 'not I', the other person, and calls this the everyday form of the character. From somebody other than me to he is the usual explanation, but the step itself is not settled.
\end{cousinweb}

\begin{grammarlens}[title={What you've built}]
One more word about people.
\end{grammarlens}

\subsection*{Your turn}
\begin{itemize}
\raggedright
  \item \emph{Say it:} \emph{tā}
  \item \emph{Say it:} \emph{tā}, once more
  \item \emph{Say it:} \emph{tā shì wǒ de érzi} — He's my son.
  \item \emph{Recall:} say the Mandarin for not bad, then the Mandarin for she, then say \emph{tā} again
\end{itemize}

\begin{tcolorbox}[breakable,colback=teal!4,colframe=teal!35!black,title={Before you move on}]
What does \zh{他} mean? (\textquotedblleft{}He; him\textquotedblright{}.) Say it once more.
\end{tcolorbox}
\section[\texorpdfstring{\zh{他} (tā)}{他 (tā)}]{\zh{他} — the person piece, then yě}

\label{lesson:ZH-W31-ta}



\begin{quote}
The word first, then the character inside it.

\begin{itemize}
\raggedright
  \item \emph{Recall:} say \emph{tā} — \textbf{R1}, one lesson back, before it decays
\end{itemize}

\begin{quote}
\zh{亻} \zh{你}
\end{quote}

\emph{Read it:} each one aloud, with nothing said first — the person piece, \emph{nǐ}
\end{quote}

\subsection*{Script — the person piece, then yě}
\hlblockfigure{\detokenize{figures/ZH-W31-ta-filmstrip.pdf}}{How \zh{他} is written, stroke by stroke}

\begin{quote}
\zh{他}
\end{quote}

\textbf{Two pieces, left before right.} Count the pieces, not the strokes.

\textbf{The left piece, \zh{亻}, two strokes.} The same two as in \zh{你}: a left-falling stroke, then a vertical down to the foot.

\textbf{The right piece, three strokes, read yě.} A horizontal that runs right, turns down and hooks. A short vertical through the middle of it. Then a long vertical that bends right along the bottom and hooks up at the end.

2 + 3 = five, and the first two your hand already knows.

\subsection*{Your turn}
\begin{itemize}
\raggedright
  \item \emph{Write it:} \zh{他}, all five, with the model in front of you
  \item \emph{Cover:} write it again from memory, and count the pieces, not the strokes
\end{itemize}

\begin{tcolorbox}[breakable,colback=teal!4,colframe=teal!35!black,title={Before you move on}]
How many strokes does \zh{他} have? (\textbf{Five}.) Which piece comes first? (\textbf{\zh{亻}, the person piece, on the left}.) How does the last stroke end? (\textbf{It runs along the bottom and hooks up}.)

Source: \href{https://raw.githubusercontent.com/chanind/hanzi-writer-data/68d10a4b21150cae5e1ebbd223eed289cf32d90c/data/\%E4\%BB\%96.json}{Hanzi Writer Data \zh{他}, pinned stroke paths and medians}.
\end{tcolorbox}
\section[\texorpdfstring{\zh{他们} (tāmen)}{他们 (tāmen)}]{\zh{他们} (tāmen) — they; them}

\label{lesson:ZH-C31-tamen}



\begin{quote}
Before the new one: say the Mandarin for she, then the Mandarin for he.
\end{quote}

\subsection*{You'll want to know: \zh{他们}}
\textbf{\zh{他们}} — \emph{tāmen} — \textquotedblleft{}they; them\textquotedblright{}.

\begin{quote}
\textbf{\zh{他们是中国人}} \emph{tāmen shì Zhōngguó rén} — They're Chinese.
\end{quote}

The new character, in pieces: \zh{们} — the person piece on the left again, and on the right a three-stroke piece read mén, the word for door, which gives the sound. The next lesson puts it in your hand.

\begin{cousinweb}
\textbf{\zh{他们}} (\emph{tāmen}) is built from \zh{他} 'he' and the plural ending men, which goes on words for people. No dictionary we use records an older form of the ending, so its history is not settled.
\end{cousinweb}

\begin{grammarlens}[title={What you've built}]
One more word about people.
\end{grammarlens}

\subsection*{Your turn}
\begin{itemize}
\raggedright
  \item \emph{Say it:} \emph{tāmen}
  \item \emph{Say it:} \emph{tāmen}, once more
  \item \emph{Say it:} \emph{tāmen shì Zhōngguó rén} — They're Chinese.
  \item \emph{Recall:} say the Mandarin for she, then the Mandarin for he, then say \emph{tāmen} again
\end{itemize}

\begin{tcolorbox}[breakable,colback=teal!4,colframe=teal!35!black,title={Before you move on}]
What does \zh{他们} mean? (\textquotedblleft{}They; them\textquotedblright{}.) Say it once more.
\end{tcolorbox}
\section[\texorpdfstring{\zh{们} (men)}{们 (men)}]{\zh{们} — the person piece, then a doorway}

\label{lesson:ZH-W31-men}



\begin{quote}
The word first, then the character inside it.

\begin{itemize}
\raggedright
  \item \emph{Recall:} say \emph{tāmen} — \textbf{R1}, one lesson back, before it decays
\end{itemize}

\begin{quote}
\zh{他} \zh{你}
\end{quote}

\emph{Read it:} each one aloud, with nothing said first — \emph{tā}, \emph{nǐ}
\end{quote}

\subsection*{Script — the person piece, then a doorway}
\hlblockfigure{\detokenize{figures/ZH-W31-men-filmstrip.pdf}}{How \zh{们} is written, stroke by stroke}

\begin{quote}
\zh{们}
\end{quote}

\textbf{Two pieces, left before right.}

\textbf{The left piece, \zh{亻}, two strokes}, exactly as in \zh{他} and \zh{你}.

\textbf{The right piece, a doorway, three strokes.} A dot at the top left. A vertical down the left side. Then a horizontal across the top that turns down the right side and hooks in at the foot.

2 + 3 = five. The doorway is read mén, and alone it is the word for door; here it gives \zh{们} its sound and loses its tone.

\subsection*{Your turn}
\begin{itemize}
\raggedright
  \item \emph{Write it:} \zh{们}, all five, with the model in front of you
  \item \emph{Cover:} write it again from memory, and count the pieces, not the strokes
\end{itemize}

\begin{tcolorbox}[breakable,colback=teal!4,colframe=teal!35!black,title={Before you move on}]
How many pieces in \zh{们}? (\textbf{Two} — \zh{亻}, then the doorway.) Which stroke of the doorway comes first? (\textbf{The dot}.) Which other character opens with the same two strokes? (\textbf{\zh{他}}, and \zh{你}.)

Source: \href{https://raw.githubusercontent.com/chanind/hanzi-writer-data/68d10a4b21150cae5e1ebbd223eed289cf32d90c/data/\%E4\%BB\%AC.json}{Hanzi Writer Data \zh{们}, pinned stroke paths and medians}.
\end{tcolorbox}
\section[\texorpdfstring{\zh{我们} (wǒmen)}{我们 (wǒmen)}]{\zh{我们} (wǒmen) — we; us}

\label{lesson:ZH-C31-women}



\begin{quote}
Before the new one: say the Mandarin for he, then the Mandarin for they.
\end{quote}

\subsection*{You'll want to know: \zh{我们}}
\textbf{\zh{我们}} — \emph{wǒmen} — \textquotedblleft{}we; us\textquotedblright{}.

\begin{quote}
\emph{wǒmen hái xíng} — We're doing all right.
\end{quote}

\begin{cousinweb}
\textbf{\zh{我们}} (\emph{wǒmen}) is built from \zh{我} 'I' and the plural ending men. \zh{我} is inherited from Middle Chinese \emph{ngaX} (Guangyun, in Baxter's transcription), which the dictionary glosses as what you call yourself; no dictionary we use records an older form of the ending, so its history is not settled.
\end{cousinweb}

\begin{grammarlens}[title={What you've built}]
One more word about people.
\end{grammarlens}

\subsection*{Your turn}
\begin{itemize}
\raggedright
  \item \emph{Say it:} \emph{wǒmen}
  \item \emph{Say it:} \emph{wǒmen}, once more
  \item \emph{Say it:} \emph{wǒmen hái xíng} — We're doing all right.
  \item \emph{Recall:} say the Mandarin for he, then the Mandarin for they, then say \emph{wǒmen} again
\end{itemize}

\begin{tcolorbox}[breakable,colback=teal!4,colframe=teal!35!black,title={Before you move on}]
What does \zh{我们} mean? (\textquotedblleft{}We; us\textquotedblright{}.) Say it once more.
\end{tcolorbox}
\section[\texorpdfstring{\zh{你们} (nǐmen)}{你们 (nǐmen)}]{\zh{你们} (nǐmen) — you all; you (more than one)}

\label{lesson:ZH-C31-nimen}



\begin{quote}
Before the new one: say the Mandarin for they, then the Mandarin for we.
\end{quote}

\subsection*{You'll want to know: \zh{你们}}
\textbf{\zh{你们}} — \emph{nǐmen} — \textquotedblleft{}you all; you (more than one)\textquotedblright{}.

\begin{quote}
\emph{nǐmen zěnmeyàng} — How are you all doing?
\end{quote}

\begin{cousinweb}
\textbf{\zh{你们}} (\emph{nǐmen}) is built from \zh{你} 'you' and the plural ending men. No dictionary we use records an older form of either, so its history is not settled.
\end{cousinweb}

\begin{grammarlens}[title={What you've built}]
One more word about people.
\end{grammarlens}

\subsection*{Your turn}
\begin{itemize}
\raggedright
  \item \emph{Say it:} \emph{nǐmen}
  \item \emph{Say it:} \emph{nǐmen}, once more
  \item \emph{Say it:} \emph{nǐmen zěnmeyàng} — How are you all doing?
  \item \emph{Recall:} say the Mandarin for they, then the Mandarin for we, then say \emph{nǐmen} again
\end{itemize}

\begin{tcolorbox}[breakable,colback=teal!4,colframe=teal!35!black,title={Before you move on}]
What does \zh{你们} mean? (\textquotedblleft{}You all; you (more than one)\textquotedblright{}.) Say it once more.
\end{tcolorbox}
